\chapter{Thuật toán và logic lõi}

Chương này trình bày đầy đủ công thức của năm thuật toán quyết định hành vi hệ thống.
Mọi hằng số đều lấy trực tiếp từ mã nguồn.

\section{Lập lịch ôn tập: FSRS với luật delta}

\subsection{Tên gọi chính xác}

Phiên bản thuật toán là \code{delta-fsrs-v3}. Tên này cố ý dài vì hai nửa của nó đến từ
hai họ thuật toán khác nhau:

\begin{itemize}[leftmargin=1.4em]
  \item \textbf{Bộ lập lịch là FSRS}: đường quên luỹ thừa nghịch đảo, độ bền
  (\en{stability}) và độ khó (\en{difficulty}) đúng theo FSRS.
  \item \textbf{Cập nhật mức nắm vững là luật delta}, \textbf{không phải} \en{Bayesian
  Knowledge Tracing}. Không có xác suất chuyển trạng thái, không có tham số trượt
  (\en{slip}) hay đoán mò (\en{guess}). Vì vậy tuyệt đối không được giả định các tính chất
  Bayes khi xây dựng tiếp trên nó.
\end{itemize}

Nhiễu do đoán mò được xử lý theo cách khác: hạ độ tin cậy của quan sát theo định dạng câu
hỏi (mục~\ref{sec:format-confidence}).

\subsection{Đường quên}

Xác suất nhớ lại sau \(t\) ngày, với độ bền \(S\) ngày:

\begin{equation}
R(t) = \left(1 + F\cdot\frac{t}{S}\right)^{D},
\qquad F = \frac{19}{81},\quad D = -\frac{1}{2}
\end{equation}

Hai hằng số này được chọn để \(R(S) = 0{,}9\). Nhờ vậy \code{stability\_days} đọc thẳng
được là \textit{"số ngày cho tới khi khả năng nhớ tụt còn 90\%"}, và khoảng ôn tập tại mục
tiêu lưu giữ mặc định đúng bằng độ bền.

Khoảng cách tới lần ôn kế tiếp là nghịch đảo chính xác của \(R\), với mục tiêu lưu giữ
\(\theta = 0{,}9\):

\begin{equation}
I(S) = \max\!\left(0{,}25,\ S\cdot\frac{\theta^{1/D} - 1}{F}\right)
\end{equation}

\begin{luuy}[Vì sao phải là nghịch đảo chính xác]
Phiên bản v1 dùng đường quên hàm mũ và một hằng số khoảng cách riêng, dẫn tới khoảng ôn
xấp xỉ \(S \times 0{,}105\). Một khái niệm phải đạt độ bền khoảng 200 ngày mới được khoảng
ôn 21 ngày, trong khi độ bền chỉ tăng \(\approx 1{,}1\) lần mỗi lần ôn — kết quả là những
lần nhớ hoàn hảo vẫn bị ghim ở mức sàn 6 giờ suốt cả chục lượt.
\end{luuy}

\subsection{Cập nhật một quan sát}

Gọi trạng thái trước là \((m, S, d)\) — mức nắm vững, độ bền, độ khó — và quan sát mới là
\((o, c)\) với \(o \in \{\text{đúng}, \text{sai}\}\), \(c \in [0,1]\) là \textbf{độ dứt
khoát} của quan sát. Đặt \(\Delta t\) là số ngày trôi qua kể từ lần tương tác trước.

\paragraph{Bước 1 — suy giảm theo thời gian.}
\begin{equation}
r = R(\Delta t),\qquad m_{\text{decay}} = \mathrm{clamp}\big(m \cdot r,\ 0{,}01,\ 0{,}99\big)
\end{equation}

\paragraph{Bước 2 — trọng số bằng chứng.} Khái niệm càng khó thì một câu trả lời càng ít
nói lên điều gì:
\begin{equation}
e = c\,(1 - 0{,}35\,d)
\end{equation}

\paragraph{Bước 3 — luật delta cho mức nắm vững.} Với \(\eta = 0{,}45\) là tốc độ học và
\(y = 1\) nếu đúng, \(y = 0\) nếu sai:
\begin{equation}
m' = \mathrm{clamp}\big(m_{\text{decay}} + e\,\eta\,(y - m_{\text{decay}}),\ 0{,}01,\ 0{,}99\big)
\end{equation}

\paragraph{Bước 4 — cập nhật độ bền.} Ba trường hợp:

\begin{equation}
S' =
\begin{cases}
S_0(o, c) & \text{lần quan sát đầu tiên}\\[4pt]
S \cdot \Big[1 + 0{,}45\,e\,(1-d)\big(1 + 1{,}6(1-r)\big)\Big] & \text{nếu đúng}\\[4pt]
\min\Big(S(1 - 0{,}55\,e),\ S_0(o,c)\Big) & \text{nếu sai}
\end{cases}
\end{equation}

Hệ số \(1{,}6\) là \textbf{thưởng giãn cách} (\en{spacing bonus}): nhớ lại được một ký ức
đã suy giảm (\(r\) nhỏ) củng cố trí nhớ mạnh hơn nhiều so với luyện lại một ký ức còn tươi.
Đây chính là hiệu ứng giãn cách mà SM-2 hoàn toàn không biểu diễn được.

Nhánh sai dùng \(\min\): chỉ nhân hệ số suy giảm là chưa đủ, vì một ký ức đã đạt 200 ngày
độ bền sau khi quên vẫn còn 90 ngày. Phải kéo hẳn về mức khởi tạo để khái niệm quay lại
giai đoạn học lại.

Độ bền khởi tạo \(S_0\) phụ thuộc chất lượng câu trả lời đầu tiên:
\begin{equation}
S_0 =
\begin{cases}
2{,}0 + 4{,}0\,c & \text{nếu đúng (2 đến 6 ngày)}\\
1{,}0 - 0{,}6\,c & \text{nếu sai (1 xuống 0,4 ngày)}
\end{cases}
\end{equation}

Lưu ý hai nhánh đi ngược chiều nhau theo cùng một đầu vào \(c\): một câu trả lời đúng dứt
khoát thì kéo dài độ bền, một câu sai dứt khoát thì rút ngắn nó.

\paragraph{Bước 5 — cập nhật độ khó.}
\begin{equation}
d' = \mathrm{clamp}\big(d \mp \delta e,\ 0,\ 1\big),
\qquad \delta = 0{,}04 \text{ (đúng)},\ 0{,}08 \text{ (sai)}
\end{equation}

Sai làm khái niệm khó lên nhanh gấp đôi so với tốc độ đúng làm nó dễ đi — thiết kế bất đối
xứng có chủ đích.

\paragraph{Bước 6 — chốt lịch.} \(S'\) được kẹp về \([0{,}25;\ 3650]\) ngày, sau đó
\(\text{next\_review\_at} = \text{now} + I(S')\).

\subsection{Quy đổi điểm 0--5 sang quan sát}

Người học chấm mức nhớ theo thang SM-2 quen thuộc (0--5), nhưng động cơ cần cặp
(kết quả, độ dứt khoát):

\begin{equation}
o = \begin{cases}\text{đúng} & q \ge 3\\ \text{sai} & q < 3\end{cases}
\qquad
c = 0{,}2 + 0{,}8 \cdot \frac{|q - 2{,}5|}{2{,}5}
\end{equation}

\begin{luuy}[Độ dứt khoát không phải là điểm]
Điểm 0 (quên sạch) và điểm 5 (nhớ ngay) đều là bằng chứng \textbf{mạnh}, chỉ khác chiều.
Điểm 2 và 3 mới là mơ hồ. Trước đây tuyến từ vựng truyền thẳng \(q/5\) làm độ tin cậy,
biến một lần quên sạch thành quan sát vô nghĩa (\(c = 0\)) và phạt câu suýt đúng nặng hơn
câu sai hoàn toàn.
\end{luuy}

\subsection{Độ tin cậy theo định dạng câu hỏi}
\label{sec:format-confidence}

Một câu trắc nghiệm bốn lựa chọn trả lời đúng có thể là đoán mò một phần tư số lần; một
câu gõ đáp án thì không. Bảng \code{FORMAT\_CONFIDENCE} nhân vào \(c\):

\begin{longtable}{L{4.4cm} L{2.6cm} L{6.6cm}}
\toprule
\textbf{Định dạng} & \textbf{Hệ số} & \textbf{Lý do}\\
\midrule
\endhead
\code{true\_false} & 0,45 & Xác suất đoán trúng 50\%\\
\code{multiple\_choice} & 0,65 & Xác suất đoán trúng \(\approx\)25\%\\
\code{matching} & 0,85 & Khó đoán nhưng vẫn có ràng buộc\\
\code{reorder} & 0,90 & Gần như không đoán được\\
\code{fill\_blank} & 1,00 & Không có gì để đoán\\
\code{translate} & 1,00 & Không có gì để đoán\\
Không rõ & 1,00 & Mặc định trung tính\\
\bottomrule
\caption{Hệ số tin cậy theo định dạng câu hỏi}
\end{longtable}

Đây chính là vai trò mà tham số \en{guess} đảm nhiệm trong BKT, đạt được ở đây bằng cách
điều chỉnh độ tin cậy mà \code{evolve\_state} vốn đã nhận vào.

\section{Chấm điểm năng lực CEFR}

\subsection{Trọng số kỹ năng}

\begin{equation}
w = \{\text{nghe}: 0{,}20,\ \text{nói}: 0{,}20,\ \text{đọc}: 0{,}20,\ \text{viết}: 0{,}20,\ \text{từ vựng}: 0{,}10,\ \text{ngữ pháp}: 0{,}10\}
\end{equation}

Cân bằng quanh bốn kỹ năng chính, mỗi kỹ năng 20\%. Từ vựng và ngữ pháp chỉ 10\% vì trong
CEFR chúng là \textbf{tài nguyên} phục vụ bốn kỹ năng, không phải trụ cột độc lập.

\subsection{Cập nhật điểm kỹ năng}

Với mỗi bài tập \(i\) thuộc kỹ năng đang xét, điểm \(s\) được cập nhật \textbf{một bước
cho mỗi bài tập}:

\begin{equation}
s \leftarrow s + \underbrace{\min\big(0{,}15,\ (1-\lambda)\cdot\mu_i\big)}_{\text{bước}} \cdot (x_i - s)
\end{equation}

trong đó \(\lambda = 0{,}95\) là hệ số giữ lịch sử, \(\mu_i\) là hệ số độ khó theo bậc
(A1: 0,5; A2: 0,7; B1: 1,0; B2: 1,3; C1: 1,6; C2: 2,0), \(x_i\) là điểm của chính bài tập
đó, và 0,15 là trần một bước.

Bốn tính chất mà công thức này bắt buộc phải giữ, mỗi tính chất tương ứng một lỗi đã từng
xảy ra:

\begin{enumerate}[leftmargin=1.4em]
  \item \textbf{Một bài tập không thể đẩy điểm lên tối đa.} Bản cũ bỏ qua phép trung bình
  khi điểm đang bằng 0, nên một câu trả lời đúng duy nhất cho ra 100/100. Nay khi kỹ năng
  chưa từng có điểm, điểm khởi đầu là \textbf{sàn của bậc hiện tại}, không phải điểm câu
  đầu tiên.
  \item \textbf{Gộp lô không đổi kết quả.} Trước đây phép cập nhật áp một lần cho mỗi lời
  gọi API. Quiz tin tức gửi theo lô còn quiz sách gửi từng câu, nên cùng sáu câu trả lời
  lại có giá trị khác nhau tới 2,4 lần. Nay áp một bước cho mỗi bài tập.
  \item \textbf{Sai không được nửa điểm.} Mục tiêu là \(x_i\) — điểm thật của bài tập —
  chứ không phải nhãn đúng/sai. Nhãn \code{is\_correct} chỉ dùng cho thống kê độ chính xác.
  \item \textbf{Độ khó phải sống sót.} Bản cũ cộng thưởng độ khó vào tử số rồi kẹp ở 100,
  khiến câu A1 và câu C2 cho kết quả y hệt. Nay độ khó nhân vào \textbf{kích thước bước}.
\end{enumerate}

\subsection{Độ tin cậy và điểm tổng hợp}

\begin{equation}
\mathrm{conf} = \min\!\left(1,\ \frac{n_{\text{tích luỹ}}}{50}\right)
\end{equation}

\(n\) là số bài tập \textbf{tích luỹ} chứ không phải số bài trong yêu cầu hiện tại. Trước
khi sửa, giá trị này đứng yên ở 0,02 mãi mãi và dòng chú thích "50 bài là đầy đủ tin cậy"
chưa từng đúng với bất kỳ người học nào.

Điểm tổng hợp — \textbf{chỉ có một công thức duy nhất}, tại \code{weighted\_overall}:

\begin{equation}
S_{\text{tổng}} = \frac{\sum_{k \in K} w_k\, s_k}{\sum_{k \in K} w_k},
\qquad K = \{\text{kỹ năng đã có điểm}\}
\end{equation}

Kỹ năng chưa đo được \textbf{loại khỏi} phép tính thay vì tính là 0, để một kỹ năng chưa
chạm tới không kéo tụt tổng điểm.

\begin{luuy}[Ba công thức từng cùng tồn tại]
Đã có lúc hệ thống có ba câu trả lời khác nhau cho "điểm tổng của người học này là bao
nhiêu": bản có trọng số (dùng để xét thăng bậc nhưng không bao giờ lưu), bản trung bình
cộng không trọng số (ghi vào hồ sơ sau mỗi bài tập), và điểm phần trăm thô từ bài kiểm tra
xếp lớp (ghi đè lên cả hai). Con số người học nhìn thấy không phải con số dùng để xét
thăng bậc của chính họ.
\end{luuy}

Ngược lại, \textbf{độ tin cậy tổng hợp} tính kỹ năng chưa đo là 0. Khác biệt này là có
chủ đích: một người chỉ luyện từ vựng thì thật sự chưa được đo đầy đủ và không nên được
thăng bậc dựa trên bằng chứng đó.

\subsection{Điều kiện thăng bậc}

Thăng bậc CEFR bị chặn bởi bốn điều kiện đồng thời, nhằm chống "cày điểm":

\begin{enumerate}[leftmargin=1.4em]
  \item Mọi kỹ năng phải đạt ngưỡng của bậc đích.
  \item Đủ số bài tập và số bài học đã hoàn thành.
  \item Đạt tỉ lệ chính xác và số ngày chuỗi tối thiểu.
  \item Độ tin cậy kỹ năng đạt \code{min\_skill\_confidence}: khối lượng trả lời câu
  \textit{"đã làm đủ chưa?"}, còn độ tin cậy trả lời câu \textit{"đã hiểu đủ về người này
  chưa?"}. Không có điều kiện thứ hai, một phiên chat dài đủ để thoả toàn bộ ngưỡng đếm.
\end{enumerate}

Ngoài ra có \textbf{cổng thi}: mỗi lần chỉ được lên tối đa một bậc, phải vượt bài thi có
bậc không thấp hơn bậc hiện tại với điểm đạt (mặc định 70\%).

\subsection{Kỹ năng đến từ nhãn, không từ phỏng đoán}

Kỹ năng mà một bài học ghi công được quyết định theo thứ tự ưu tiên:
\code{Lesson.skill} $\rightarrow$ \code{Course.skill} $\rightarrow$ suy đoán từ thẻ
(\code{infer\_skill\_from\_tags}). Nhánh cuối chỉ là di sản: nó mặc định trả về
\code{VOCABULARY}, và đó chính là lý do nội dung nghe/nói từng bị ghi công sai kỹ năng.

Trò chơi ánh xạ qua bảng cố định \code{GAME\_TYPE\_SKILL}, không bao giờ tách chuỗi
\code{game\_type}:

\begin{longtable}{L{4.2cm} L{3.4cm} L{6.0cm}}
\toprule
\textbf{Loại trò chơi} & \textbf{Kỹ năng} & \textbf{Ghi chú}\\
\midrule
\endhead
\code{word\_scramble} & Từ vựng & \\
\code{matching} & Từ vựng & \\
\code{spelling\_bee} & Từ vựng & \\
\code{hangman} & Từ vựng & \\
\code{fill\_blank} & \textbf{Ngữ pháp} & Rút câu từ ngân hàng ngữ pháp; trước đây bị tính nhầm sang từ vựng\\
\code{grammar\_quiz} & Ngữ pháp & \\
\bottomrule
\caption{Ánh xạ trò chơi sang kỹ năng CEFR}
\end{longtable}

\section{Bộ nhớ đệm SCAR-L1}

\subsection{Hai lớp đệm}

\begin{itemize}[leftmargin=1.4em]
  \item \textbf{L0} — khớp chính xác theo dấu vân tay (\en{fingerprint}) gồm: câu hỏi đã
  chuẩn hoá, ý định, bậc CEFR, tiếng mẹ đẻ, tập khái niệm gốc và số lượt hội thoại. Thời
  gian sống 3.600 giây.
  \item \textbf{L1} — khớp gần theo \en{bucket} đồ thị: các yêu cầu khác chữ nhưng cùng
  cụm khái niệm và cùng trạng thái. Thời gian sống 7.200 giây.
\end{itemize}

\subsection{Điểm rủi ro tái sử dụng}

Với mỗi ứng viên trong L1, hệ thống tính điểm rủi ro \(\rho \in [0,1]\):

\begin{equation}
\rho = 0{,}20\,\delta_{\text{intent}}
     + 0{,}30\,(1 - J_c)
     + 0{,}15\,(1 - J_r)
     + 0{,}20\,\delta_{\text{evidence}}
     + 0{,}15\,\max(a,\ \delta_{\text{version}})
\end{equation}

trong đó:

\begin{longtable}{L{3.0cm} L{10.6cm}}
\toprule
\textbf{Ký hiệu} & \textbf{Ý nghĩa}\\
\midrule
\endhead
\(\delta_{\text{intent}}\) & 0 nếu cùng ý định, 1 nếu khác.\\
\(J_c\) & Hệ số Jaccard giữa hai tập khái niệm (giao chia hợp).\\
\(J_r\) & Hệ số Jaccard giữa hai tập gợi ý quan hệ.\\
\(\delta_{\text{evidence}}\) & 0 nếu cùng băm bằng chứng, 1 nếu khác.\\
\(a\) & Tỉ lệ tuổi của khoản mục so với thời gian sống, kẹp trong \([0,1]\).\\
\(\delta_{\text{version}}\) & 1 nếu bất kỳ phiên bản nào lệch: \en{epoch} hồ sơ, phiên bản chính sách, phiên bản đồ thị, phiên bản nguồn, lớp độ tươi.\\
\bottomrule
\caption{Thành phần điểm rủi ro SCAR-L1}
\end{longtable}

\subsection{Quyết định ba nhánh}

\begin{equation}
\text{quyết định} =
\begin{cases}
\texttt{full} & \text{nếu vi phạm bất kỳ ràng buộc cứng nào}\\
\texttt{reuse} & \text{nếu câu hỏi chuẩn hoá trùng khớp và } \rho \le \tau_0 = 0{,}25\\
\texttt{patch} & \text{nếu } \rho \le \tau_1 = 0{,}55\\
\texttt{full} & \text{ngược lại}
\end{cases}
\end{equation}

\textbf{Ràng buộc cứng} được kiểm tra \textbf{trước} khi tính rủi ro, và bất kỳ vi phạm nào
cũng buộc tính lại toàn bộ:

\begin{itemize}[leftmargin=1.4em]
  \item Phiên bản chứng chỉ phải đúng (\code{CERTIFICATE\_SCHEMA\_VERSION} = 3).
  \item Mọi chiều bắt buộc phải có mặt ở cả hai phía. Chiều bắt buộc cơ sở gồm: ý định,
  bậc, \en{epoch} hồ sơ, phiên bản chính sách, phiên bản đồ thị, mục tiêu câu trả lời.
  \item Bậc CEFR, tiếng mẹ đẻ, \en{epoch} hồ sơ, ý định và mục tiêu câu trả lời phải khớp
  \textbf{tuyệt đối}.
  \item Băm bằng chứng, phiên bản chính sách/đồ thị/nguồn, lớp độ tươi phải khớp nếu có.
  \item Độ chồng lấp khái niệm \(J_c\) phải \(\ge 0{,}50\).
  \item Khoản mục không được quá hạn.
\end{itemize}

\begin{luuy}[Vì sao khoá đệm không chứa mức nắm vững]
Nếu bỏ điểm nắm vững vào khoá đệm, mỗi người học sẽ có một không gian đệm riêng và tỉ lệ
trúng đệm gần như bằng không. Giải pháp là dùng \code{state\_epoch} — một số nguyên duy
nhất tăng khi trạng thái đổi. Khoản mục cá nhân hoá gắn phạm vi người dùng bằng HMAC cộng
\en{epoch}; khoản mục đồ thị con dùng chung thì \textbf{không bao giờ} chứa mức nắm vững
hay định danh người dùng thô.
\end{luuy}

\subsection{Vô hiệu hoá}

Bộ đệm bị vô hiệu theo bộ bốn phiên bản
\(\langle \nu_{\text{graph}}, \nu_{\text{policy}}, \nu_{\text{profile}}, t_{\text{refresh}} \rangle\):
tăng \(\nu_{\text{graph}}\) huỷ mọi \en{bucket} L1; tăng \(\nu_{\text{policy}}\) huỷ cả L0
lẫn L1; tăng \(\nu_{\text{profile}}\) chỉ huỷ khoản mục của người dùng đó; hết hạn
\(t_{\text{refresh}}\) huỷ theo thời gian. Hệ thống còn duy trì \textbf{cạnh ngược}
(\en{reverse edges}) từ phụ thuộc sang tạo tác, để khi một nguồn bằng chứng đổi thì mọi
phản hồi đã suy ra từ nó đều bị gỡ.

\section{Hoà trộn truy xuất}

\subsection{Điểm hợp nhất}

Mỗi mẩu bằng chứng nhận một điểm hợp nhất từ ba tín hiệu:

\begin{equation}
s = \alpha\,\underbrace{\frac{1}{1+h}}_{\text{đồ thị}} + \beta\,\underbrace{\cos(q, e)}_{\text{véc-tơ}} + \gamma\,\underbrace{e^{-\lambda \Delta t}}_{\text{độ mới}}
\end{equation}

với \(\alpha = 0{,}5\), \(\beta = 0{,}3\), \(\gamma = 0{,}2\), \(\lambda = 0{,}01\);
\(h\) là số bước nhảy trên đồ thị tới khái niệm gốc, \(\cos(q,e)\) là độ tương đồng cô-sin
từ MiniLM, \(\Delta t\) là số lượt hội thoại kể từ lần dùng gần nhất.

Trọng số nghiêng về cấu trúc đồ thị (0,5) hơn tương đồng ngữ nghĩa (0,3), vì trong dạy
ngôn ngữ thì quan hệ tiên quyết đáng tin hơn độ giống nhau bề mặt của câu chữ.

\subsection{Bộ xếp hạng học trực tuyến}

Sau hợp nhất, danh sách được xếp lại bằng bộ xếp hạng lô-gít theo cặp
(\en{pairwise logistic ranker}) có cập nhật trực tuyến nhẹ. Mười hai đặc trưng và trọng số
khởi tạo:

\begin{longtable}{L{4.4cm} L{2.2cm} L{7.0cm}}
\toprule
\textbf{Đặc trưng} & \textbf{Trọng số} & \textbf{Ý nghĩa}\\
\midrule
\endhead
\code{base\_score} & 0,42 & Điểm hợp nhất ở bước trước\\
\code{body\_overlap} & 0,14 & Chồng lấp từ với thân đoạn\\
\code{title\_overlap} & 0,11 & Chồng lấp từ với tiêu đề\\
\code{anchor\_coverage} & 0,11 & Độ phủ các thực thể neo trong câu hỏi\\
\code{title\_phrase} & 0,08 & Khớp cụm nguyên vẹn ở tiêu đề\\
\code{title\_token\_recall} & 0,07 & Độ bao phủ từ khoá của tiêu đề\\
\code{kg\_proximity} & 0,05 & Khoảng cách trên đồ thị tri thức\\
\code{vec\_signal} & 0,05 & Tín hiệu tương đồng véc-tơ\\
\code{graph\_signal} & 0,05 & Tín hiệu cấu trúc đồ thị\\
\code{memory\_signal} & 0,04 & Từng dùng trong phiên này\\
\code{recency\_signal} & 0,03 & Độ mới\\
\code{external\_penalty} & \(-0{,}05\) & Phạt nguồn ngoài đồ thị\\
\bottomrule
\caption{Đặc trưng của bộ xếp hạng truy xuất}
\end{longtable}

Trọng số được cập nhật trực tuyến với chuẩn hoá \(L_2\) và kẹp trong \([-3;\ 3]\) để một
đợt phản hồi bất thường không thể lật đổ toàn bộ mô hình.

\subsection{Chỉ số đánh giá}

\code{EvaluationAgent} tính hai nhóm chỉ số. Nhóm chất lượng văn bản: WER
(\en{Word Error Rate} — tỉ lệ lỗi từ, bằng khoảng cách soạn thảo chia độ dài), TTR
(\en{Type-Token Ratio} — số từ khác nhau chia tổng số từ, đo độ phong phú từ vựng) và mật
độ lỗi. Nhóm chất lượng truy xuất:

\begin{equation}
\mathrm{P@}k = \frac{|\text{liên quan} \cap \text{lấy ra}|}{k},\quad
\mathrm{R@}k = \frac{|\text{liên quan} \cap \text{lấy ra}|}{|\text{liên quan}|}
\end{equation}
\begin{equation}
\mathrm{NDCG@}k = \frac{\sum_{i=1}^{k} \frac{rel_i}{\log_2(i+1)}}{\mathrm{IDCG@}k},
\qquad
\mathrm{MRR} = \frac{1}{\text{hạng của mẩu liên quan đầu tiên}}
\end{equation}

\section{Chống cày điểm trong gamification}

Hệ thống XP có bốn tầng bảo vệ:

\begin{longtable}{L{4.6cm} L{9.0cm}}
\toprule
\textbf{Cơ chế} & \textbf{Chi tiết}\\
\midrule
\endhead
Trần ngày & \code{DAILY\_XP\_CAP} = 500 XP mỗi ngày cho toàn bộ nguồn.\\
Trần một lần cấp & \code{MAX\_SINGLE\_AWARD} = 200 XP; vượt quá thì từ chối yêu cầu.\\
Trần theo nguồn & Trò chơi 100, bài học 50, thử thách ngày 50, sách 25, podcast 20, tin tức 15, YouTube 15.\\
Bắt buộc \code{source\_id} & Ba nguồn nhạy cảm lặp — \code{game}, \code{lesson}, \code{daily\_challenge} — bắt buộc kèm \code{source\_id} do máy chủ cấp. Thiếu thì trả lỗi 422. Chỉ mục duy nhất một phần trong cơ sở dữ liệu chặn cấp trùng.\\
\bottomrule
\caption{Cơ chế chống cày XP}
\end{longtable}

Hệ số chuỗi ngày nhân vào XP cơ sở: chuỗi \(\ge\)3 ngày nhân 1,2; \(\ge\)7 ngày nhân 1,5;
\(\ge\)30 ngày nhân 2,0.

\section{Ranh giới ghi nhận: cái gì được tính, cái gì không}

Đây là quy tắc dễ vi phạm nhất khi thêm tính năng mới:

\begin{longtable}{L{5.0cm} L{4.2cm} L{4.4cm}}
\toprule
\textbf{Hoạt động} & \textbf{Điểm kỹ năng} & \textbf{Trạng thái khái niệm}\\
\midrule
\endhead
Hoàn thành bài học & Có & Có (nếu bài tập có \code{concept\_id})\\
Trò chơi & Có & Có\\
Ôn từ vựng & Có & Có (\textbf{không} truyền lại \code{concept\_id}, tránh đếm hai lần)\\
Quiz tin tức / sách & Có & Có\\
Lượt chat với AI & Có (qua kênh quan sát) & Có\\
Nghe podcast, xem YouTube, đọc thuần tuý & \textbf{Không} & \textbf{Không}\\
\bottomrule
\caption{Hoạt động nào tạo bằng chứng học tập}
\end{longtable}

\begin{luuy}[Vì sao tiêu thụ thụ động không ghi nhận gì]
Nghe hết một tập podcast không phải bằng chứng của việc hiểu. Các bề mặt này sẽ được nối
vào hệ thống ghi nhận khi chúng có bài chính tả hoặc quiz kèm theo — nghĩa là khi có thứ
để chấm.
\end{luuy}

Ngoài ra, tham số \code{award\_xp=False} phải được truyền bất cứ khi nào tuyến gọi đã tự
cấp XP (bài học, trò chơi, ôn từ, lượt chat), vì XP trong hàm ghi nhận là \textbf{cộng
thêm} chứ không thay thế.
